% ============================================================================
% REQUIRED PACKAGES - Essential for document rendering
% ============================================================================

% Mathematical typesetting (required for equations and symbols)
\usepackage{amsmath,amssymb}          % Mathematical symbols and environments
\usepackage{amsfonts}                 % Additional math fonts
\usepackage{amsthm}                   % Theorem environments

% Graphics and page layout (required for figures and formatting)
\usepackage{graphicx}                 % Include graphics (REQUIRED for figures)
\usepackage{geometry}                 % Page margins (set via pandoc --variable)
\usepackage{float}                    % Better float placement

% Tables (required for table formatting)
\usepackage{booktabs}                 % Professional tables
\usepackage{longtable}                % Long tables spanning pages
\usepackage{array}                    % Advanced table formatting

% PDF features (required for cross-references and metadata)
\usepackage{url}                      % URL formatting
\usepackage{hyperref}                 % Hyperlinks and cross-references
% natbib is loaded once by the Pandoc template; options are applied in _render_pdf_override._postprocess_tex

% ============================================================================
% PACKAGES - Improve formatting and functionality
% ============================================================================

% Table enhancements (optional but recommended)
\usepackage{multirow}                 % Multi-row table cells
\usepackage{caption}                  % caption formatting
\usepackage{subcaption}               % Sub-figures and sub-tables

% Math enhancements (optional but recommended)
\usepackage{bm}                       % Bold math symbols

% Reference enhancements (optional but recommended)
\usepackage{cleveref}                 % Intelligent cross-referencing
% Render DOIs as inline hyperlinks (not footnotes).  The doi.sty package
% is intentionally omitted because it renders DOIs as footnotes.  This
% override typesets them as clickable links within the bibliography entry.
\providecommand{\doi}[1]{\href{https://doi.org/#1}{\nolinkurl{doi:#1}}}

% Configure figure numbering and captions
\renewcommand{\figurename}{Figure}
\captionsetup{
    justification=centering,
    font=small,
    labelfont=bf,
    labelsep=period
}

% Configure table numbering and captions
\renewcommand{\tablename}{Table}
\captionsetup[table]{
    justification=centering,
    font=small,
    labelfont=bf,
    labelsep=period
}

% Configure section numbering
\setcounter{secnumdepth}{3}
\renewcommand{\thesection}{\arabic{section}}
\renewcommand{\thesubsection}{\arabic{section}.\arabic{subsection}}
\renewcommand{\thesubsubsection}{\arabic{section}.\arabic{subsection}.\arabic{subsubsection}}

% Configure equation numbering
\numberwithin{equation}{section}

% Configure hyperref for proper linking
\hypersetup{
    colorlinks=true,
    linkcolor=red,
    citecolor=red,
    urlcolor=red,
    filecolor=red,
    anchorcolor=red,
    pdfborder={0 0 0},
    bookmarks=true,
    bookmarksnumbered=true,
    bookmarkstype=toc,
    pdftitle={Ento-Linguistics: Language, Ambiguity, and Scientific Communication in Entomology},
    pdfsubject={Scientific Terminology Analysis in Entomology},
    pdfproducer={XeLaTeX}
    % pdfauthor and pdfkeywords are set dynamically in _title_page.tex
    % from config.yaml to avoid metadata drift
}

% ============================================================================
% PACKAGE CONFIGURATION
% ============================================================================

% Configure cleveref for intelligent cross-references
\crefname{section}{Section}{Sections}
\crefname{subsection}{Subsection}{Subsections}
\crefname{subsubsection}{Subsubsection}{Subsubsections}
\crefname{equation}{Equation}{Equations}
\crefname{figure}{Figure}{Figures}
\crefname{table}{Table}{Tables}
\crefname{appendix}{Appendix}{Appendices}

% Configure fonts for Unicode support with fallbacks
\usepackage{newunicodechar}
%\newunicodechar definitions for Unicode characters used in the manuscript
\newunicodechar{⁴}{\textsuperscript{4}}
\newunicodechar{₄}{\textsubscript{4}}
\newunicodechar{²}{\textsuperscript{2}}
\newunicodechar{₀}{\textsubscript{0}}
\newunicodechar{₁}{\textsubscript{1}}
\newunicodechar{₂}{\textsubscript{2}}
\newunicodechar{₃}{\textsubscript{3}}
% Arrow characters used in abstract, conclusion, and supplemental methods
\newunicodechar{→}{\ensuremath{\rightarrow}}
\newunicodechar{←}{\ensuremath{\leftarrow}}
\newunicodechar{↔}{\ensuremath{\leftrightarrow}}
% Additional symbols
\newunicodechar{≥}{\ensuremath{\geq}}
\newunicodechar{≤}{\ensuremath{\leq}}
\newunicodechar{≠}{\ensuremath{\neq}}
\newunicodechar{…}{\ldots}

% ============================================================================
% FONTS AND TYPOGRAPHY
% ============================================================================

% Use standard fonts for better compatibility
% NOTE: T1 fontenc is intentionally omitted — it is incompatible with xelatex
% (which uses Unicode natively). Including it generates font warnings that
% can cascade into compilation failures.
\usepackage{lmodern}

% ============================================================================
% CODE BLOCK STYLING
% ============================================================================

% code block styling for better contrast and readability
\usepackage{fancyvrb}
\usepackage{fvextra}   % Extends fancyvrb: adds bgcolor, breaklines, etc.
\usepackage{xcolor}
\usepackage{listings}

% Define custom colors for code blocks
\definecolor{codebg}{RGB}{248, 248, 248}      % Very light gray background
\definecolor{codeborder}{RGB}{200, 200, 200}  % Medium gray border
\definecolor{codefg}{RGB}{34, 34, 34}         % Dark gray text
\definecolor{commentcolor}{RGB}{102, 102, 102} % Comment color
\definecolor{keywordcolor}{RGB}{0, 0, 0}       % Keyword color
\definecolor{stringcolor}{RGB}{0, 102, 0}      % String color

% Configure Verbatim environment for inline code
\DefineVerbatimEnvironment{Verbatim}{Verbatim}{%
    fontsize=\small,
    frame=single,
    framerule=0.5pt,
    framesep=3pt,
    rulecolor=\color{codeborder},
    bgcolor=codebg
}

% Configure code block styling
\DefineVerbatimEnvironment{Highlighting}{Verbatim}{%
    fontsize=\footnotesize,
    frame=single,
    framerule=0.5pt,
    framesep=5pt,
    rulecolor=\color{codeborder},
    bgcolor=codebg
}

% Style inline code with \texttt
\renewcommand{\texttt}[1]{%
    \colorbox{codebg}{\color{codefg}\ttfamily #1}%
}

% Configure listings package for code blocks
\lstset{
    backgroundcolor=\color{codebg},
    basicstyle=\footnotesize\ttfamily\color{codefg},
    breakatwhitespace=false,
    breaklines=true,
    captionpos=b,
    commentstyle=\color{commentcolor},
    deletekeywords={...},
    escapeinside={\%*}{*)},
    extendedchars=true,
    frame=single,
    framerule=0.5pt,
    framesep=5pt,
    keepspaces=true,
    keywordstyle=\color{keywordcolor}\bfseries,
    language=Python,
    morekeywords={*,...},
    numbers=left,
    numbersep=5pt,
    numberstyle=\tiny\color{codefg},
    rulecolor=\color{codeborder},
    showspaces=false,
    showstringspaces=false,
    showtabs=false,
    stepnumber=1,
    stringstyle=\color{stringcolor},
    tabsize=4,
    title=\lstname
}

% Override any Pandoc default lstset configurations
\AtBeginDocument{
    \lstset{
        backgroundcolor=\color{codebg},
        basicstyle=\footnotesize\ttfamily\color{codefg},
        frame=single,
        framerule=0.5pt,
        framesep=5pt,
        rulecolor=\color{codeborder},
        numbers=left,
        numbersep=5pt,
        numberstyle=\tiny\color{codefg}
    }
}

% Configure bibliography
% Note: Using plainnat with natbib package for proper citation processing
% The bibliography style and commands (\bibliographystyle and \bibliography) are in 99_references.md

% Simple page break support for document structure
% Note: Page breaks are handled in the markdown generation, not here

% ============================================================================
% DOCUMENT FORMATTING
% ============================================================================

% Ensure proper spacing and formatting
\frenchspacing  % Single space after periods
\linespread{1.2}  % Slightly increased line spacing for readability

% ============================================================================
% TITLE PAGE OVERRIDE
% ============================================================================
% The title page, author block, DOI, and PDF metadata are generated at
% build time by _render_pdf_override.py from config.yaml, written to
% _title_page.tex, and \input'd here.  This avoids hardcoding metadata
% in two places.  The generated file redefines \maketitle using \parbox
% blocks instead of the default tabular (which breaks under TeXLive 2026).
\IfFileExists{_title_page.tex}{% Auto-generated from config.yaml — do not edit by hand.
% Uses a custom \maketitle to avoid TeXLive 2026 tabular/\and breakage.
\makeatletter
\renewcommand{\maketitle}{%
  \begin{center}
    {\LARGE \@title \par}
    \vskip 1.5em
    \parbox[t]{0.45\textwidth}{\centering
      \textbf{Daniel Ari Friedman}\\
      {\small Active Inference Institute, APOIDEAS}\\
      {\footnotesize {\ttfamily daniel@activeinference.institute} \enspace $\cdot$ \enspace ORCID: 0000-0001-6232-9096}
    }
    \hfill
    \parbox[t]{0.45\textwidth}{\centering
      \textbf{Tucker Cahill Chambers}\\
      {\small APOIDEAS}\\
      {\footnotesize ORCID: 0009-0008-3793-7872}
    }
    \vskip 1em
    April 15, 2026\\[4pt]
    {\small DOI: \href{https://doi.org/10.5281/zenodo.19574118}{10.5281/zenodo.19574118}}
  \end{center}
  \par\vskip 1.5em
}
\makeatother

% PDF metadata (synced from config.yaml keywords and authors)
\hypersetup{%
  pdfauthor={Daniel Ari Friedman and Tucker Cahill Chambers},
  pdfkeywords={entomology, scientific terminology, language analysis, discourse analysis, ant biology, eusociality, conceptual frameworks, active inference, terminology networks, Markov blankets, stigmergy, computational linguistics},
}
}{}
